\section{Modélisation prédictive du coût opérationnel}
\label{modelisation}
\shorthandoff{;:!?}

Ce chapitre exploite la base de mesure construite au chapitre~\ref{construction-mesure} pour répondre aux deux finalités du mémoire~: reconstruire le coût rétrospectif d'un deal traité, et prédire le coût d'un deal futur avant son traitement. Fidèle à la dualité établie (Event Processing et Aftercare), il construit deux modèles distincts, chacun adapté à la nature de sa composante, puis les recombine en un coût total prédit.

\subsection{Reconstruire le coût rétrospectif d'un deal}
\label{sec:cout-retro}

La première exploitation est directe~: appliquer l'équation du TDABC à la base de mesures empiriques. Pour un deal donné, le coût opérationnel est la somme, sur tous ses événements, du temps standard valorisé au coût horaire chargé~:

\begin{equation}
C_{deal} = \sum_{i=1}^{N} \frac{T_{s,i}}{60}\,w_s
= \underbrace{\frac{N\,\bar{t}}{60}\,w_s}_{\text{Event Processing}}
\;+\; \underbrace{\frac{T_{AC}}{60}\,w_s}_{\text{Aftercare}}
\end{equation}

\noindent où $N$ est le volume d'événements du deal, $\bar{t}\approx 20$~minutes est le temps unitaire empirique d'Event Processing, $T_{\text{AC}}$ le temps d'Aftercare, et $w_{s}$ le coût horaire chargé du site traitant (ici estimé à 60~€/h en moyenne).
La force de cette reconstruction est qu'elle rend visible ce qui était opaque~: on peut désormais dire non seulement qu'un deal «~coûte cher~», mais \emph{pourquoi} --- combien de minutes viennent de l'Event Processing, combien de l'Aftercare, quelle part est imputable au volume et quelle part à la structure.

% Préambule requis : \usetikzlibrary{positioning, arrows.meta, calc}  + \usepackage{booktabs}
\begin{figure}[htbp]
\centering
\begin{tikzpicture}[
  >=Stealth, font=\small,
  panel/.style={rectangle, rounded corners=7pt, draw=#1!55, fill=#1!5,
                line width=1pt, inner sep=9pt},
  res/.style={rectangle, rounded corners=6pt, draw=#1!70!black, fill=#1!14,
              line width=1pt, align=center, font=\footnotesize\bfseries,
              minimum width=30mm, minimum height=11mm},
  kpi/.style={rectangle, rounded corners=6pt, draw=#1!70!black, fill=#1!12,
              line width=1pt, align=center, font=\scriptsize,
              minimum width=22mm, minimum height=14mm},
  op/.style={font=\normalsize\bfseries, text=gray!70},
  flow/.style={->, line width=1pt, color=gray!55},
]

% ============ PANNEAU EVENT PROCESSING ============
\node[panel=blue] (ep) {%
  \begin{tabular}{@{}l r r r@{}}
    \multicolumn{4}{@{}l}{\textcolor{blue!55!black}{\bfseries Event Processing}\;\footnotesize\textcolor{gray}{— booking Loan~IQ (mesuré)}}\\[3pt]
    \toprule
    \textbf{Événement} & \textbf{N} & \textbf{$\bar t$} & \textbf{min}\\
    \midrule
    Repricings     & 2\,855 & 22 & 62\,810\\
    Transfers      & 1\,928 & 20 & 38\,560\\
    Drawdowns      & 1\,399 & 25 & 34\,975\\
    Fees           &   871  & 18 & 15\,678\\
    Paper clip     &   621  & 19 & 11\,799\\
    Accrual cycle  &   679  & 16 & 10\,864\\
    Adjustments    &   205  & 14 &  2\,870\\
    Others         &   104  & 20 &  2\,080\\
    \midrule
    \textbf{Total} & \textbf{8\,662} & \textbf{$\approx$21} & \textbf{179\,636}\\
    \bottomrule
  \end{tabular}%
};

% ============ PANNEAU AFTERCARE ============
\node[panel=orange, anchor=north west] (ac) at ([xshift=12mm]ep.north east) {%
  \begin{tabular}{@{}l r@{}}
    \multicolumn{2}{@{}l}{\textcolor{orange!60!black}{\bfseries Aftercare (SAV)}}\\
    \multicolumn{2}{@{}l}{\footnotesize\textcolor{gray}{post-traitement (estimé)}}\\[4pt]
    \toprule
    Rôle              & Agent \\
    Prêteurs          & 600 \\
    Facilités         & 8 \\
    Événements aval   & 31\,820 \\
    Mails (Hobart)    & 50\,000 \\
    \midrule
    \textbf{$\hat T_{\text{AC}}$} & \textbf{$\approx$79\,000 min} \\
    \bottomrule
  \end{tabular}%
};

% ============ NŒUD TOTAL (convergence) ============
\coordinate (mid) at ($(ep.south)!0.5!(ac.south)$);
\node[res=teal, below=16mm of mid] (tot) {Charge totale\\$\approx$ 259\,000 min};

\draw[flow] (ep.south) |- ([yshift=6mm]tot.north) -- (tot.north);
\draw[flow] (ac.south) |- ([yshift=6mm]tot.north);
\node[op] at ($(ep.south)!0.5!(ac.south)!0.55!(tot.north)$) {$+$};

% ============ CHAÎNE DE CONVERSION ============
\node[kpi=teal,   below=14mm of tot] (h)   {\textbf{$\approx$4\,317}\\heures};
\node[kpi=blue,   below=10mm of h]   (etp) {\textbf{$\approx$2,7}\\ETP};
\node[kpi=red,    below=10mm of etp] (eur) {\textbf{$\approx$259\,000}\\EUR};

\draw[flow] (tot) -- (h)   node[midway, above, font=\scriptsize, text=gray]{$\div\,60$};
\draw[flow] (h)   -- (etp) node[midway, above, font=\scriptsize, text=gray]{$\div\,1\,600$};
\draw[flow] (etp) -- (eur) node[midway, above, font=\scriptsize, text=gray]{$\times\,60\,\text{EUR}/h$};

% ============ FORMULE ============
\node[below=7mm of tot.south west, anchor=north west, font=\footnotesize]
  {$C = (T_{\text{LIQ}} + T_{\text{AC}}) \times N \times w_s
     \;\approx\; 30\ \text{min/évt} \times 8\,662 \times 60\,\text{EUR}/h$};

\end{tikzpicture}
\caption{Exemple de calcul du coût d'un deal (rôle Agent, profil \emph{jumbo}, anonymisé). Le \textit{booking} est mesuré par clocking ($\bar t \approx 20$~min/événement, 8\,662 événements) ; l'\textit{aftercare} est estimé à partir des inducteurs du deal. Les deux composantes forment une charge totale de $\approx$~259\,000~minutes, soit $\approx$~2,7~ETP, valorisée à $\approx$~259\,000~\EUR{} au coût horaire chargé de 60~\EUR/h.}
\label{fig:exemple-cout-deal}
\end{figure}


\paragraph{Un résultat structurant~: la dispersion extrême des coûts.}
Appliquée au portefeuille, cette reconstruction révèle une dispersion considérable du coût par deal. Un deal simple --- rôle de participant, peu de prêteurs, monodevise, faible volume --- se traite pour l'équivalent de quelques milliers d'euros. Un deal complexe --- rôle d'agent, plusieurs centaines de prêteurs, multidevise, volume d'événements élevé --- peut atteindre plusieurs centaines de milliers d'euros.

Ainsi, un \textbf{deal standard} (participant, 2 prêteurs, 1 facilité, monodevise, $\approx$~198 événements) se traite pour $\approx$~\textbf{4\,440~€} (0,05~ETP), tandis qu'un \textbf{deal jumbo} (agent, $\approx$~600 prêteurs, 8 facilités, multidevise, $\approx$~8\,662 événements) peut atteindre $\approx$~\textbf{259~k€} (2,7~ETP).

\begin{figure}[H]
    \centering
    \begin{tikzpicture}[
        >=stealth',
        profile/.style={draw=blue!80!black, rectangle, rounded corners, align=center,
                        fill=blue!10, font=\small, text width=5cm, minimum height=1.4cm},
        cost/.style={draw=blue!80!black, thick, rectangle, rounded corners, align=center,
                     fill=blue!15, font=\small, text width=5cm, minimum height=1.4cm},
        synth/.style={draw=green!60!black, thick, rectangle, rounded corners, align=center,
                      fill=green!20, font=\small, text width=11cm, minimum height=1.2cm}
    ]
        % --- Profils : deux colonnes comparées ---
        \node[profile] (simple)
            {\textbf{Deal Standard} (Participant) \\
             {\footnotesize 2 prêteurs, 1 facilité, monodevise, OCR oui} \\
             {\footnotesize Volume faible — 198 events}};
        \node[profile, right=2.5cm of simple] (complex)
            {\textbf{Jumbo Deal} (Agent) \\
             {\footnotesize $>1000$ prêteurs, 4 facilités, multidevise, OCR non} \\
             {\footnotesize Volume élevé — 8 662 events}};

        % --- Coûts respectifs, alignés sous chaque profil ---
        \node[cost, below=1.2cm of simple] (cost-simple)
            {\textbf{Coût opérationnel bas} \\
             {\footnotesize $\approx 21$ min/event $\rightarrow$ \textbf{4 440 €} (0,05 FTE)}};
        \node[cost, below=1.2cm of complex] (cost-complex)
            {\textbf{Coût opérationnel très élevé} \\
             {\footnotesize $\approx 30$ min/event $\rightarrow$ \textbf{259 000 €} (2,7 FTE)}};

        \draw[->, thick] (simple) -- (cost-simple);
        \draw[->, thick] (complex) -- (cost-complex);

        % --- Synthèse centrée sous les deux coûts ---
        \node[synth, below=1.6cm of {$(cost-simple)!0.5!(cost-complex)$}] (ratio)
            {\textbf{Dispersion extrême des coûts — rapport de 1 à 60} \\
             {\footnotesize 1 \% des deals génèrent 30 \% du volume global d'événements Back Office}};

        \draw[->, dashed, thick] (cost-simple.south) -- (ratio.north);
        \draw[->, dashed, thick] (cost-complex.south) -- (ratio.north);
    \end{tikzpicture}
    \caption{Impact de la complexité structurale sur le coût opérationnel et la charge Back Office}
    \label{fig:dispersion-couts-style}
\end{figure}



\paragraph{Ce que la mesure débloque.}
En rendant ce coût comparable au gain connu du Front Office, la reconstruction rétrospective répond au constat fondateur~: contrairement au Front Office, dont l'activité se mesure en revenus (P\&L), le Back Office ne se mesurait qu'en coût global, jamais par deal. Le coût par deal ouvre trois usages immédiats~: chiffrer l'impact des initiatives de réduction des coûts (objectifs GTS), comparer le coût de traitement entre sites pour évaluer le \textit{cost avoidance} du nearshoring, et identifier les deals à coût anormalement élevé.

\subsection{Identifier les facteurs explicatifs de la charge}
\label{sec:facteurs}

Reconstruire le coût ne suffit pas~: pour le \emph{prédire} et l'\emph{expliquer}, il faut savoir quels paramètres le déterminent. La séparation des deux composantes impose de distinguer deux ensembles de facteurs, de nature différente.

\paragraph{Les facteurs de l'Event Processing~: qualitatifs et catégoriels.}
Le volume d'événements $N$ dépend de variables essentiellement \emph{qualitatives} définies dès la signature du deal~:
\begin{itemize}
    \item le \textbf{rôle} de la banque (agent, participant, bilatéral)~;
    \item le \textbf{nombre de prêteurs} (\textit{lenders}) impliqués~;
    \item le \textbf{nombre de facilités} (\textit{facilities})~;
    \item le \textbf{produit} et la \textbf{ligne de métier}~.
\end{itemize}
Ces facteurs agissent de façon \emph{non additive}~: l'effet du rôle d'agent n'est pas le même selon le nombre de prêteurs, et se combine avec lui. C'est cette structure d'interactions qui oriente le choix de modèles à base d'arbres (Random Forest, XGBoost) pour prédire $\hat{N}(x)$, le volume d'événements attendu.



\paragraph{Les facteurs de l'Aftercare~: quantitatifs et proportionnels.}
Le temps d'Aftercare, lui, dépend de variables \emph{quantitatives} qui mesurent le volume de gestion aval~:
\begin{itemize}
    \item le \textbf{nombre de prêteurs} ($n_{\text{lenders}}$)~;
    \item le \textbf{nombre de facilités} ($n_{\text{facilities}}$)~;
    \item le \textbf{nombre de courriels} de gestion ($n_{\text{mails}}$)~;
    \item le \textbf{nombre d'événements} ($N$)~.
\end{itemize}
Leur relation à la charge est approximativement \emph{proportionnelle}~: plus un deal compte de prêteurs, plus il génère de documents, de confirmations et de réconciliations. Cette proportionnalité rend la relation modélisable par une régression (linéaire ou polynomiale) aux coefficients directement interprétables --- «~chaque prêteur supplémentaire ajoute $X$ minutes d'Aftercare~» --- langage immédiatement actionnable pour le management.

\subsection{Construire un modèle prédictif du coût}
\label{sec:modele-predictif}

La seconde finalité --- prédire le coût d'un deal \emph{avant} son traitement --- consiste à apprendre, sur l'historique, la fonction qui relie les caractéristiques d'un deal à sa charge. Fidèle à la dualité, nous construisons \textbf{deux modèles distincts}, puis les recombinons.

\begin{figure}[htbp]
\centering
\begin{tikzpicture}[
    >=stealth',
    box/.style={rectangle, draw, fill=blue!10, rounded corners,
                text width=4cm, align=center, minimum height=1.1cm, font=\small},
    modelbox/.style={rectangle, draw=blue!80!black, thick, fill=blue!15, rounded corners,
                     text width=4cm, align=center, minimum height=1.7cm, font=\small},
    result/.style={rectangle, draw=green!60!black, thick, fill=green!20, rounded corners,
                   text width=11cm, align=center, minimum height=1.3cm, font=\small}
]

    % Deux étages côte à côte (écart suffisant pour ne pas se chevaucher)
    \node[modelbox] (N)
        {\textbf{Étage 1 — Volume $\hat{N}$}\\ {\footnotesize Arbres / XGBoost}\\ $\hat{N}(x)\times 20\,\text{min}$};
    \node[modelbox, right=3cm of N] (AC)
        {\textbf{Étage 2 — Aftercare $\hat{T}_{\text{AC}}$}\\ {\footnotesize Régression paramétrique}\\ $\hat{T}_{\text{AC}}(x)$};

    % Entrée centrée au-dessus des deux étages
    \node[box, above=1.6cm of {$(N)!0.5!(AC)$}] (input)
        {\textbf{Caractéristiques du deal}\\ {\footnotesize Rôle, prêteurs, facilités, multidevise, produit}};

    % Résultat centré en dessous
    \node[result, below=1.8cm of {$(N)!0.5!(AC)$}] (final)
        {\textbf{Coût total prédit $\hat{C}_{\text{deal}}$}\\[2pt]
         $\hat{C}_{\text{deal}} = \bigl[\,\hat{N}(x)\times 20\,\text{min} + \hat{T}_{\text{AC}}(x)\,\bigr]\times 60\,\text{€/h}$};

    % Connecteurs
    \draw[->, thick] (input) -- (N);
    \draw[->, thick] (input) -- (AC);
    \draw[->, thick] (N) -- (final);
    \draw[->, thick] (AC) -- (final);
\end{tikzpicture}
\caption{Architecture à deux étages du modèle prédictif du coût par deal}
\label{fig:modele-deux-etages}
\end{figure}



\subsubsection{Modéliser le volume d'événements~: Random Forest et XGBoost}
\label{sec:modele-ep}

Les facteurs de volume d'événements étant qualitatifs et en interaction, une régression linéaire échouerait --- son hypothèse d'additivité ne tient pas. Nous retenons donc des \textbf{méthodes à base d'arbres}, qui capturent nativement les interactions.

\begin{itemize}
    \item Le \textbf{Random Forest} moyenne des centaines d'arbres entraînés sur des échantillons aléatoires : $\hat{N}(x) = \frac{1}{B}\sum_{b=1}^{B} T_b(x)$. Robuste, peu sensible au réglage, il fournit une mesure d'importance des variables qui hiérarchise les facteurs de volume.
    \item Le \textbf{gradient boosting} (XGBoost) construit les arbres séquentiellement, chacun corrigeant les erreurs du précédent~; généralement plus précis sur données tabulaires, au prix d'un réglage plus fin.
\end{itemize}

L'interprétabilité perdue par l'agrégation est restaurée \emph{a posteriori} par les valeurs SHAP, qui décomposent chaque prédiction en contributions additives --- permettant de présenter au management, pour un deal donné, la part de volume imputable à chaque facteur structurel.


\begin{figure}[htbp]\centering
\begin{tikzpicture}[>=Stealth, font=\scriptsize,
  dec/.style={draw, rounded corners, fill=blue!8, align=center, text width=24mm, inner sep=3pt},
  leaf/.style={draw, rounded corners, fill=green!12, align=center, text width=18mm, inner sep=3pt},
  el/.style={font=\scriptsize\itshape, fill=white, inner sep=1pt}]
  \node[dec] (r) {Rôle = Agent ?};
  \node[leaf, below left=12mm and 22mm of r] (p) {Participant\\$\approx$ 4 k€};
  \node[dec, below right=12mm and 6mm of r] (n1) {Prêteurs $>$ 50 ?};
  \node[leaf, below left=12mm and 0mm of n1] (l1) {$\approx$ 40 k€};
  \node[dec, below right=12mm and 0mm of n1] (n2) {Multidevise ?};
  \node[leaf, below left=10mm and -2mm of n2] (l2) {$\approx$ 120 k€};
  \node[leaf, below right=10mm and -2mm of n2] (l3) {$\approx$ 259 k€};
  \draw[->] (r) -- node[el,left]{non} (p);
  \draw[->] (r) -- node[el,right]{oui} (n1);
  \draw[->] (n1) -- node[el,left]{non} (l1);
  \draw[->] (n1) -- node[el,right]{oui} (n2);
  \draw[->] (n2) -- node[el,left]{non} (l2);
  \draw[->] (n2) -- node[el,right]{oui} (l3);
\end{tikzpicture}
\caption{Exemple illustratif d'arbre de décision~: chaque chemin de la racine à une feuille combine des caractéristiques du deal (rôle, prêteurs, multidevise) et aboutit à un coût estimé. L'arbre capture ainsi les effets d'interaction que la régression linéaire manque. \textbf{Valeurs illustratives}}
\label{fig:arbre-illustratif}
\end{figure}

\begin{table}[H]
\centering
\caption{Performance du modèle de volume d'événements (Random Forest vs XGBoost)}
\label{tab:metriques-performance}
\begin{tabular}{lccc}
\toprule
\textbf{Modèle} & \textbf{$R^2$} & \textbf{MAE (min)} & \textbf{RMSE (min)} \\
\midrule
Random Forest   & 0.65 & 45.2 & 67.8 \\
XGBoost        & 0.8 & 38.5 & 52.3 \\
\bottomrule
\end{tabular}
\small
\vspace{2mm}
\\
\textit{Note : Évaluation sur un échantillon de 7000 deals traités en 2025.}
\end{table}

\medskip
\paragraph{Ce que le modèle estime, et quand.}
Le modèle ne prédit pas directement un coût, mais le \textbf{nombre d'événements $\hat{N}(x)$} qu'un deal générera sur sa vie opérationnelle, à partir des seules caractéristiques \emph{connues dès le closing}~: rôle de la banque, nombre de facilités, produit, multidevise et nombre de prêteurs initial. C'est ce volume, une fois multiplié par le temps unitaire empirique ($\bar t \approx 20$~min) et valorisé au coût horaire chargé, qui donne l'estimation de coût. L'intérêt métier est là~: on obtient une première évaluation de la charge \emph{avant même d'avoir traité le moindre événement}, au moment où les arbitrages --- dimensionnement d'équipe, priorité d'automatisation, choix de site --- se décident.

\medskip
\paragraph{Une estimation vivante, non un verdict.}
Certaines caractéristiques d'entrée sont figées dès la signature (le rôle, le produit), d'autres sont amenées à évoluer. Le nombre de prêteurs, en particulier, se stabilise au fil de la syndication~: un deal signé avec un tour de table indicatif peut voir ce nombre croître ou se recomposer avant la mise en place. La prédiction faite au closing est donc une \textbf{borne initiale}, destinée à être rafraîchie à mesure que la structure du deal se fige. Loin d'être une faiblesse, cette propriété fait du modèle un outil de pilotage \emph{dynamique}~: on ne fournit pas un coût définitif, mais une estimation qui se resserre à chaque jalon, et dont l'écart avec la volumétrie réellement observée devient lui-même un signal (un deal qui «~déborde~» sa prédiction initiale mérite attention).

\medskip
\paragraph{Une prédiction qui se lit comme une règle métier.}
Le choix d'arbres n'est pas seulement technique~: il préserve la lisibilité. L'arbre de la figure~\ref{fig:arbre-illustratif} se parcourt comme une suite de questions qu'un opérateur poserait spontanément --- \textit{la banque est-elle Agent~? le tour de table dépasse-t-il 50 prêteurs~? le deal est-il multidevise~?} Chaque chemin de la racine à une feuille est une règle explicite, et la progression des feuilles matérialise le régime de charge correspondant. On ne présente pas au management une boîte noire, mais une partition intelligible du portefeuille, directement traduisible en langage opérationnel.

\medskip
\paragraph{Hiérarchiser les inducteurs de charge.}
Au-delà d'un deal isolé, le Random Forest fournit une \textbf{mesure d'importance des variables}, calculée à partir de la réduction de variance cumulée qu'apporte chaque facteur. Elle répond à une question que le management se pose directement~: \emph{qu'est-ce qui, structurellement, fait le volume~?} Si le rôle et le nombre de prêteurs dominent, l'action prioritaire n'est pas d'optimiser le traitement unitaire, mais d'anticiper les deals \textit{Agent} à large syndication. Cette hiérarchie oriente l'allocation de l'effort d'automatisation là où le gain marginal est le plus fort.

\medskip
\paragraph{Décomposer le coût, deal par deal, face au Front Office.}
La précision gagnée par l'agrégation d'arbres se paie en lisibilité~: on ne lit plus une règle, seulement un score. Les valeurs SHAP la restaurent \emph{a posteriori}, en décomposant chaque prédiction en contributions \emph{additives} qui somment exactement à l'écart au coût moyen. On peut alors tenir au Front Office un discours chiffré et défendable~: «~ce deal est estimé à 180~k€, soit 90~k€ au-dessus de la moyenne~; $+50$~k€ viennent du nombre de prêteurs, $+30$~k€ du rôle d'Agent, $+10$~k€ du multidevise~». C'est la condition pour qu'une estimation de coût opposée à une équipe commerciale soit argumentable sur des données plutôt que ressentie --- exigence non négociable en contexte bancaire.


\subsubsection{Modéliser l'Aftercare~: régression pour l'interprétabilité}
\label{sec:modele-ac}

Pour l'Aftercare, le choix s'inverse. La relation étant proportionnelle et les variables quantitatives, une \textbf{régression} est non seulement suffisante mais préférable. Le modèle de base s'écrit~:

\begin{equation}
\hat{T}_{\text{AC}} = \beta_0 + \beta_1 \, n_{\text{lenders}} + \beta_2 \, N
+ \beta_3 \, n_{\text{mails}} + \beta_4 \, n_{\text{facilities}} + \varepsilon
\label{eq:reg-aftercare}
\end{equation}

\noindent où chaque coefficient $\beta_k$ se lit directement comme le temps d'Aftercare ajouté par une unité supplémentaire du facteur correspondant. Là où la relation présente une courbure --- par exemple un effet d'échelle au-delà d'un certain nombre de prêteurs --- une \textbf{régression polynomiale} l'absorbe sans sacrifier l'interprétabilité.
En prédiction, $N$ n'étant pas encore observé, on y substitue le volume prédit $\hat N(x)$ issu de l'étage~1~; l'erreur sur $\hat N$ se propage donc dans $\hat T_{\text{AC}}$. Ces régresseurs sont par ailleurs corrélés entre eux (plus de prêteurs implique mécaniquement plus d'événements et de courriels)~: cette colinéarité rend les coefficients individuels moins stables, et l'interprétation «~effet marginal d'un facteur~» doit être maniée avec prudence.




\subsubsection{Recomposer le coût prédit}
\label{sec:cout-predit-recompo}

Les deux modèles se recombinent en un coût total prédit pour un deal caractérisé par son vecteur d'attributs $x$~:

\begin{equation}
c_u = \frac{w_s}{60} = \frac{60\ \text{€/h}}{60} = 1\ \text{€/min}
\qquad\Longrightarrow\qquad
\hat{C}_{deal} = \big[\hat{N}(x)\times 20 + \hat{T}_{AC}(x)\big]\times c_u
\label{eq:cout-predit-total}
\end{equation}

\noindent où $\hat{N}(x)$ est le volume d'événements prédit par le modèle à base d'arbres, $\hat{T}_{\text{AC}}(x)$ est le temps d'Aftercare prédit par la régression, et le facteur 60~€/h est le coût horaire chargé moyen estimé.

Cette équation est le produit central du mémoire du côté prédictif~: elle permet, dès la signature d'un deal --- lorsqu'on connaît son rôle, son nombre de prêteurs, de facilités et de produit, sans avoir traité le moindre événement --- d'estimer ce qu'il coûtera aux opérations. La performance s'évalue par les métriques usuelles (MAE, RMSE, $R^2$), en distinguant l'erreur sur chaque composante pour diagnostiquer laquelle limite la précision globale.

\subsection{Pilotage opérationnel et usages décisionnels}
\label{sec:pilotage}

La double mesure (rétrospective et prédictive) alimente quatre leviers de pilotage concrets.

\paragraph{1. Cibler les initiatives de réduction des coûts.}
Puisque 1\% des deals génèrent 30\% du volume d'événements, la modélisation permet d'identifier \emph{ex ante} les deals à haut volume probable et de concentrer les efforts d'automatisation ou de simplification processuelle sur cette frange critique. Le gain potentiel est disproportionné~: réduire de 10\% le volume d'événements sur les deals du top 1\% a plus d'impact que la même réduction sur l'ensemble du portefeuille.

\paragraph{2. Évaluer le \textit{cost avoidance} du nearshoring.}
Le coût par deal étant désormais comparable entre sites (Paris / Lisbonne), on peut simuler le coût d'un même portefeuille traité dans l'une ou l'autre configuration, à structure de deals identique. La différence observée, nette de l'effet volume, quantifie le \textit{cost avoidance} réel du transfert d'activité.

\paragraph{3. Anticiper la charge et dimensionner les équipes.}
Le modèle prédictif, appliqué au pipeline de deals signés mais pas encore traités, fournit une projection de charge opérationnelle (en ETP et en euros) sur les mois à venir. Cela transforme la planification d'une logique budgétaire statique en une logique dynamique fondée sur la volumétrie attendue.
Symétriquement, cette projection fournit un argument objectif pour \emph{défendre}
un dimensionnement~: en cas de pression sur les effectifs, une surcharge réelle peut
être démontrée chiffres à l'appui, plutôt que ressentie.

\paragraph{4. Mesurer le gain des initiatives d'automatisation et d'IA.} 
Parce que la charge est désormais chiffrée avant et après, on peut évaluer l'impact réel d'une
initiative --- RPA, OCR (DocFactory), IA générative --- en euros et en ETP économisés
sur les familles d'événements concernées. La mesure devient l'instrument commun de
justification des initiatives stratégiques de la banque~: elle transforme un « gain
estimé » en gain constaté, deal par deal et famille par famille.


\subsection{Limites et prolongements}
\label{sec:limites}

\paragraph{Biais résiduels et contextualisation.}
La mesure porte sur des opérateurs conscients d'être observés, dans un contexte de nearshoring où la mesure peut être perçue comme liée à la maîtrise des coûts. L'effet Hawthorne n'a pas été éliminé, seulement contenu et rendu explicite. Les temps mesurés sont donc des ordres de grandeur robustes, non des chronométrages absolus.

\paragraph{Variables non observées.}
Certaines caractéristiques influençant la charge ne sont pas dans les données disponibles~: complexité juridique du contrat, qualité de la documentation initiale, présence de clauses atypiques, criticité réputationnelle du deal. Ces facteurs contribuent au bruit résiduel des modèles.

\paragraph{Extrapolation et stabilité temporelle.}
Le modèle prédit pour des deals futurs dont la distribution peut diverger de l'historique (nouveaux produits, changements réglementaires, évolution des pratiques de marché). Une gouvernance de ré-entraînement périodique et de surveillance de la \textit{distribution shift} est nécessaire pour maintenir la validité prédictive.

\paragraph{Prolongements possibles.}
Les directions naturelles d'extension sont~: intégrer les données de Hobart (e-mails) en variables explicatives plus fines pour l'Aftercare~; explorer des modèles hybrides combinant régression et arbres pour capturer conjointement proportionnalité et interactions~; étendre la prédiction à l'horizon multi-mois pour la planification capacitaire~; et tester la transférabilité du modèle à d'autres métiers du Back Office (Trade Finance, Securities Services).

\shorthandon{;:!?}